\NativeSlot{45}{67}{%
% p030_s01

\ExamQuestion{6（续）}{5}%
\(182\times66=(\square+1)\times66=181\times66+1\times\square\)。\par

% p030_s02

\noindent 所以\(181\times67\ \bigcirc\ 182\times66\)。发现：当两个数之和相同时，如果两个数之间的差越小，那么这两个数的积就越\mbox{（\hspace{1.8em}）}。\par
}
\NativeSlot{161}{36}{%
% p030_s03

\ExamQuestion{（2）}{5}%
根据发现来比较：\(9999\times2026\ \bigcirc\ 9998\times2027\)。\par
}
\NativeSlot{188}{39}{%
% p030_s04

\ExamQuestion{7.}{5}%
\((81+36)\div9=81\div9+36\div9\)，\allowbreak{}\((121+22)\div11=121\div11+22\div11\)，请你照样子写一组这样的算式\mbox{（\hspace{9em}）}。\par

% p030_s05

\noindent 用字母表示上面的规律\mbox{（\hspace{9em}）}。\par
}
\NativeSlot{264}{43}{%
\ExamGroup{二、选择（将正确答案的序号填在括号里）。{\normalfont\kaishu\fontsize{10.8bp}{14bp}\selectfont （每题2分，共12分）}}

% p030_s06

\ExamQuestion{1.}{5}%
\(53\times99=53\times100-53\times1\)应用了\mbox{（\hspace{1.8em}）}。\par
\ExamOptions{A. 加法交换律}{B. 乘法交换律}{C. 乘法结合律}{D. 乘法分配律}
}
\NativeSlot{339}{43}{%
% p030_s07

\ExamQuestion{2.}{5}%
下面\mbox{（\hspace{1.8em}）}的算式和\(27\times99\)的结果相等。\par
\ExamOptions{A. \(27\times90\times9\)}{B. \(27\times100-1\)}{C. \(27\times100-27\)}{D. \(27\times100-100\)}
}
\NativeSlot{415}{36}{%
% p030_s08

\ExamQuestion{3.}{5}%
如果\(a\times(b\times18)=216\)，那么\(a\times b=\)\mbox{（\hspace{1.8em}）}。\par
\ExamOptions{A. 17}{B. 198}{C. 3888}{D. 12}
}
\NativeSlot{463.0}{43}{%
% p030_s09

\ExamQuestion{4.}{5}%
小虎在一次计算中把\(5\times(a+6)\)错写成了\(5\times(a+9)\)，计算结果比原来\mbox{（\hspace{1.8em}）}。\par
\ExamOptions{A. 增加15}{B. 减少15}{C. 增加3}{D. 减少3}
}
\NativeSlot{538}{74}{%
\par\noindent%
\ExamTile{.56}{%
% p030_s10

\ExamQuestion{5.}{5}%
如图，能用“\((15+24)\times32\)”解决的问题是\mbox{（\hspace{1.8em}）}。\par
\ExamOptions{A. 种西红柿、白菜和豆角的面积一共多少平方米}{B. 种白菜和豆角的面积一共多少平方米}{C. 种西红柿和豆角的面积一共多少平方米}{D. 种西红柿和白菜的面积一共多少平方米}
}\hfill%
\ExamTile{.415}{%
\BookFigure{figures/p030_s10-01.tex}
}\par
}
\NativeSlot{662.5}{78}{%
% p030_s11

\ExamQuestion{6.}{5}%
如图，能说明“\(6\times3+4\times3\)”与“\((6+4)\times3\)”相等的是\mbox{（\hspace{1.8em}）}。\par
\ExamOptions{A. ①②}{B. ②③}{C. ①②③}{D. ①②③④}
\par\noindent%
\ExamTile{.485}{%
\BookFigure{figures/p030_s11-01.tex}
}\hfill%
\ExamTile{.485}{%
\BookFigure{figures/p030_s11-02.tex}
}\par
}
